\documentclass{article}
\usepackage[T1]{fontenc}
\usepackage{ifthen}
\usepackage{tikz}
\usetikzlibrary{math}
\usepackage[left=1.5cm,right=1.5cm,top=2cm,bottom=2cm]{geometry}
\usepackage{url}
\usepackage{minted}
\providecommand{\mintex}[1]{\mintinline[breaklines]{tex}{#1}}

% Open the file for reading
\newread\myfile
% \def\size{6}
% \openin\myfile=example.txt
\def\size{100}
\openin\myfile=input.txt

\title{Advent of code : 2015 day 18}
\author{In \LaTeX\ }
\date{Many years later !}


\def\setpos#1#2#3{\expandafter\xdef\csname LIGHT-#1-#2-\endcsname{#3}}
\def\getpos#1#2{\csname LIGHT-#1-#2-\endcsname}
\def\setmpos#1#2#3{\expandafter\def\csname TMP-#1-#2-\endcsname{#3}}
\def\getmpos#1#2{\csname TMP-#1-#2-\endcsname}

\newcount\nbon
\newcount\around
\newcount\ip
\newcount\im
\newcount\jp
\newcount\jm

\newif\ifcontinue
\makeatletter

% Fill the current line of the main array, from the input string
\def\treatLine#1{%
  \expandafter\@treatLine#1\relax
}
\def\@treatLine#1{%
    \advance\idj by1
    \ifx\relax#1%
        % End of the string
    \else
        \ifx #1.%
            % Off
            % \the\idi\ \the\idj\ is Off
            \else
            % On
            % \the\idi\ \the\idj\ is On
            \setpos{\the\idi}{\the\idj}{1}
        \fi
        \expandafter\@treatLine
    \fi
}

% Display the array of light in a tikz figure
% Also displays the number of light that are on
\def\display{%
    \nbon=0
    \begin{tikzpicture}[scale=0.1]
        \draw (0, 0) -- (\size+1, 0) -- (\size+1, \size+1) -- (0, \size+1) -- cycle;
        \foreach \i in {1,2,...,\size}{%
            \foreach \j in {1,2,...,\size}{%
                \if\getpos{\i}{\j}1
                    \global\advance\nbon by1
                    \draw[fill=magenta] (\i, \j) circle[radius=0.5cm];
                % \else
                    % \draw[] (\i, \j) circle[radius=0.5cm];
                \fi
            }
        }
        \draw (\size+1, \size) node[anchor=north west] {\the\nbon\ are On};
    \end{tikzpicture}
}

\def\counton{%
    \nbon=0%
    \foreach \i in {1,2,...,\size}{%
        \foreach \j in {1,2,...,\size}{%
            \if\getpos{\i}{\j}1%
                \global\advance\nbon by1%
                % F\,\the\nbon
            \fi%
        }%
    }%
    \the\nbon
}

%Copy the content of the temp array into the array
\def\copyback{%
    \foreach \i in {1,2,...,\size}{%
        \foreach \j in {1,2,...,\size}{%
            \global\setpos{\i}{\j}{\getmpos{\i}{\j}}
        }
    }
}


\def\lightstep{%
    \foreach \i in {1,2,...,\size}{%
        \foreach \j in {1,2,...,\size}{%
            \around=0

            \ip=\i
            \advance\ip by1
            \im=\i
            \advance\im by-1
            \jp=\j
            \advance\jp by1
            \jm=\j
            \advance\jm by-1

            % \i, \j\ :
            \if\getpos{\the\im}{\the\jm}1
                \advance\around by1
            \fi
            \if\getpos{\the\im}{\j}1
                \advance\around by1
            \fi
            \if\getpos{\the\im}{\the\jp}1
                \advance\around by1
            \fi

            \if\getpos{\i}{\the\jm}1
                \advance\around by1
            \fi
            \if\getpos{\i}{\the\jp}1
                \advance\around by1
            \fi

            \if\getpos{\the\ip}{\the\jm}1
                \advance\around by1
            \fi
            \if\getpos{\the\ip}{\j}1
                \advance\around by1
            \fi
            \if\getpos{\the\ip}{\the\jp}1
                \advance\around by1
            \fi

            % \phantom{www}\the\around%
            \if\getpos{\i}{\j}1
                \if\the\around 2
                    \global\setmpos{\i}{\j}{1}
                    % Stays on (2)
                \else\if\the\around 3
                    \global\setmpos{\i}{\j}{1}
                    % Stays on (3)
                \else
                    \global\setmpos{\i}{\j}{0}
                    % Get off
                \fi
                \fi
            \else
                \if\the\around 3
                    \global\setmpos{\i}{\j}{1}
                    % Get ON !
                \else
                    \global\setmpos{\i}{\j}{0}
                    % remains OFF
                \fi
            \fi
        }
    }
}

\begin{document}

\maketitle

\section{Read the input file}

%%%%%%%%%%%%%%%%%%%%%%%%%%%%%%%%%%%%%%%%% Text 🡓

In this document, we define a « two dimensional » array to store the lights\footnote{Source: \url{https://tex.stackexchange.com/a/367490/243474}.}.
It defines two macros \mintex{\setpos} and \mintex{\getpos} that allow to respectively set and read the value at a given row and column in the array.

\begin{minted}{tex}
\def\setpos#1#2#3{\expandafter\xdef\csname LIGHT-#1-#2-\endcsname{#3}}
\def\getpos#1#2{\csname LIGHT-#1-#2-\endcsname}
\end{minted}

To fill this array, we read the input file and treat each line one after the other.
To do so, we use two cunts \mintex{\idi} and \mintex{\idj} that will be the indices of the light in the array, and read the file line by line into \mintex{\content}
Then for each line, we use a separate macro, \mintex{\treatLine{\content}} that will recursively treat each character of the line.
If the character is a \texttt{.}, then nothing happens, and if it's another character (in the present case, it is always a \texttt{\#}) then the corresponding case of the array is set to \texttt{1}.
In these process, we manually increment the values of \mintex{\idi} and \mintex{\idj}.

We also define a macro \mintex{\size} which is the total size of the grid (\texttt{100} in the input, but \texttt{6} in the example).

\begin{minted}{tex}
% Open the file for reading
\newread\myfile
\openin\myfile=input.txt

\newcount\idi
\idi=0
\newcount\idj
\idj=0

% Define a macro to store the contents of the file
\newtoks\content

\continuetrue
\loop
\read\myfile to \content

\ifeof\myfile
    \continuefalse
\else
    \advance\idi by1
    \idj=0
    % \content
    \treatLine{\content}
    % \parseline{\content}
    % Line is \content
\fi
\ifcontinue
\repeat

% Close the file
\closein\myfile
\end{minted}

\begin{minted}{tex}
\newif\ifcontinue
\makeatletter
% Fill the current line of the main array, from the input string
\def\treatLine#1{%
  \expandafter\@treatLine#1\relax
}
\def\@treatLine#1{%
    \advance\idj by1
    \ifx\relax#1%
        % End of the string
    \else
        \ifx #1.%    % Off
        \else%       % On
            \setpos{\the\idi}{\the\idj}{1}
        \fi
        \expandafter\@treatLine
    \fi
}
\end{minted}



%%%%%%%%%%%%%%%%%%%%%%%%%%%%%%%%%%%%%%%%% Code 🡓


\newcount\idi
\idi=0
\newcount\idj
\idj=0

% Define a macro to store the contents of the file
\newtoks\content

\continuetrue
\loop
\read\myfile to \content

\ifeof\myfile
    \continuefalse
\else
    \advance\idi by1
    \idj=0
    % \content
    \treatLine{\content}
    % \parseline{\content}
    % Line is \content
\fi
\ifcontinue
\repeat

% Close the file
\closein\myfile



%%%%%%%%%%%%%%%%%%%%%%%%%%%%%%%%%%%%%%%%% Text 🡓

\section{Display the table, and count the lights}

By making a nested loop, we can display the array and count the light that are on.
Precisely, we use a count initially at 0 and use two \mintex{foreach} loops.
Note that we need to use the \mintex{\global} keywords inside the loop\footnote{\url{https://tex.stackexchange.com/a/766457/243474}.}.

We define two macros \mintex{\display} and \mintex{\counton} that do that.
Note the second is identical to the first on, but without the display using Ti\emph{k}Z (so I do not rewrite the code here).

\begin{minted}{tex}
\newcount\nbon
\def\display{%
    \nbon=0
    \begin{tikzpicture}[scale=0.1]
        \draw (0, 0) -- (\size+1, 0) -- (\size+1, \size+1) -- (0, \size+1) -- cycle;
        \foreach \i in {1,2,...,\size}{%
            \foreach \j in {1,2,...,\size}{%
                \if\getpos{\i}{\j}1
                    \global\advance\nbon by1
                    \draw[fill=magenta] (\i, \j) circle[radius=0.5cm];
                \else
                    % To avoid excedding TeX capacity, we only draw on lights
                    % \draw[] (\i, \j) circle[radius=0.5cm];
                \fi
            }
        }
        \draw (\size+1, \size) node[anchor=north west] {\the\nbon\ are On};
    \end{tikzpicture}
}
\end{minted}

Tiny remark : because of default axis orientation, there is a quarter-turn rotation in what is displayed here compared to what is on advent of code page. From here, apply a rotation of $-\frac\pi2$ to get what is displayed on the problem page.


%%%%%%%%%%%%%%%%%%%%%%%%%%%%%%%%%%%%%%%%% Code 🡓

\counton


\begin{center}
    \display
\end{center}






%%%%%%%%%%%%%%%%%%%%%%%%%%%%%%%%%%%%%%%%% Text 🡓

\section{Before getting in the hard, some utils macros}

In the following, we will apply the rule for switching on and off the lights.
As to compute light at time step $k$ we use the state at time $k-1$ that cannot be changed, we need to get a temporary array on which we will set the results, and in the end the values of this temporary array are applied back to the main light array.

First we define the macros \mintex{\setmpos} and \mintex{\getmpos} that are identical as the one for the light array, and we define the \mintex{\copyback} macro :

\begin{minted}{tex}
\def\copyback{%
    \foreach \i in {1,2,...,\size}{%
        \foreach \j in {1,2,...,\size}{%
            \global\setpos{\i}{\j}{\getmpos{\i}{\j}}
        }
    }
}
\end{minted}

%%%%%%%%%%%%%%%%%%%%%%%%%%%%%%%%%%%%%%%%% Code 🡓

%% just for tests

% \foreach\i in {0,1,...,100}{
%     \foreach\j in {0, 2, 13}
%     {
%         \global\setmpos{\i}{\j}{1}
%     }
% }


% \copyback

% \display
% \counton



%%%%%%%%%%%%%%%%%%%%%%%%%%%%%%%%%%%%%%%%% Text 🡓

\section{One step}

To simulate one step of switching on and off the lights, we just apply the algorithm described.
This may not be the best way to do it, but I use define one count that count the number of light that are on.
And other counts that contain the values of $i\pm1$ and $j\pm1$.
These counters need to be defined outside the macro, this is due to \TeX\ behavior, that will define a new one without erasing the previous ones, and eventually gets out of room to create new count.

\begin{minted}{tex}
\newcount\around
\newcount\ip
\newcount\im
\newcount\jp
\newcount\jm

\def\lightstep{%
    \foreach \i in {1,2,...,\size}{%
        \foreach \j in {1,2,...,\size}{%
            \around=0
            \ip=\i
            \advance\ip by1
            \im=\i
            \advance\im by-1
            \jp=\j
            \advance\jp by1
            \jm=\j
            \advance\jm by-1

            \if\getpos{\the\im}{\the\jm}1   % i-1, j-1
                \advance\around by1
            \fi
            \if\getpos{\the\im}{\j}1        % i-1, j
                \advance\around by1
            \fi
            \if\getpos{\the\im}{\the\jp}1   % i-1, j+1
                \advance\around by1
            \fi
            \if\getpos{\i}{\the\jm}1        % i, j-1
                \advance\around by1
            \fi
            \if\getpos{\i}{\the\jp}1        % i, j+1
                \advance\around by1
            \fi
            \if\getpos{\the\ip}{\the\jm}1   % i+1, j-1
                \advance\around by1
            \fi
            \if\getpos{\the\ip}{\j}1        % i+1, j
                \advance\around by1
            \fi
            \if\getpos{\the\ip}{\the\jp}1   % i+1, j+1
                \advance\around by1
            \fi

            % \phantom{www}\the\around%
            \if\getpos{\i}{\j}1
                \if\the\around 2
                    \global\setmpos{\i}{\j}{1}
                    % Stays on (2)
                \else\if\the\around 3
                    \global\setmpos{\i}{\j}{1}
                    % Stays on (3)
                \else
                    \global\setmpos{\i}{\j}{0}
                    % Get off
                \fi
                \fi
            \else
                \if\the\around 3
                    \global\setmpos{\i}{\j}{1}
                    % Get ON !
                \else
                    \global\setmpos{\i}{\j}{0}
                    % remains OFF
                \fi
            \fi
        }
    }
}
\end{minted}

Then we just go on !

\begin{minted}{tex}
\lightstep
\copyback
\end{minted}

%%%%%%%%%%%%%%%%%%%%%%%%%%%%%%%%%%%%%%%%% Code 🡓





\display


\section{100 steps}

We just make a loop to call \mintex{\lightstep} and \mintex{\copyback} the desired number of time.
As right above we already made one iteration, we begin the loop at 2.
I also comment the line to display, because with the huge grid the document would be too large (and it might reach \TeX\ capacity...)

\begin{minted}{tex}
\UseName{int_step_inline:nnn}{2}{100}{%
    \lightstep
    \copyback
    Iteration #1: \counton
    % \display
    \par
}
\end{minted}



\UseName{int_step_inline:nnn}{2}{100}{%
    \lightstep
    \copyback
    Iteration #1: \counton
    % \display
    \par
}

\display

\begin{center}
    Solution of part I : \counton
\end{center}

\end{document}
